% page 351 du fac-simile (image p-350 du scan)
% bloc 172.7 x 242.0 mm ; marge gauche 29.3 mm
\pgc{7.620mm}{0.000mm}{
\l{\cel{55}343}
\l{}
\l{\sou{locionar}. (\sou{trans.}) Fluigar liquido sur parto di la homo-korpo}
\l{\cel{3}(o sur la korpo tota), sive per varsar olu, sive per ekir-}
\l{\cel{3}igar olu per preso ek sponjo, linjo humidigita, e c. - DEFIS.}
\l{}
\l{\sou{logo}. (\sou{maro-navig.}) Planketo balastizita, imersita ye la extrem-}
\l{\cel{3}ajo di kordo, ye la dopa parto di navo, por mezurar la rapi-}
\l{\cel{3}deso-grado di la iro da la navo. - DEFIR.}
\l{}
\l{\sou{logaritmo}. (\sou{matem.}) Exponento integra o fracionala di la po-}
\l{\cel{3}tenco a qua oportas elevar nombro konstanta por reproduktar}
\l{\cel{3}ta o ca nombro. - DEFIRS.}
\l{}
\l{\sou{logiko}. I. La cienco pri lo exakta. - II. La cienco pri la +leyi}
\l{\cel{3}(legi) qui koncernas rezonado. - III. La maniero rezonar}
\l{\cel{3}rigoroze. - IV. Interligeso racionala di la kozi. (ex.:\cel{1}\sou{che}}
\l{\cel{3}\sou{Ido}\cel{1}: la arto kompozar la frazi o la vorti tale ke la adi-}
\l{\cel{3}ciono di la senco di singla ek la elementi donas (sen arbi-}
\l{\cel{3}trialeso, sen deviaco de la normo) la senco preciza di la}
\l{\cel{3}ensemblo di la frazo o di la vorto). - DEFIRS.}
\l{}
\l{\sou{logogrifo}. Sorto di enigmato qua konsistas ye un vorto di qua}
\l{\cel{3}la literi, kombinita diverse, formacas altra vorti qui anke}
\l{\cel{3}esas divinenda. - DEFIRS.}
\l{}
\l{\sou{logomakio}. Disputo pri vorti. - DEFIRS.}
\l{}
\l{\sou{lojar}. (\sou{netrans., en, che})\cel{2}Esar instalita sub tekto, en domo}
\l{\cel{3}od en parto di domo, sive kontinue, sive dum-kelka-tempe.DEFIRS.}
\l{}
\l{\sou{lojio}. I. Chambreto aden qua onu enklozas la dementi. - II. Mikra}
\l{\cel{3}domo di domo-gardisto. - III. Chambreto aden qua onu enklozas}
\l{\cel{3}la studenti pri estetiko (piktado, skultado, muzikado, e c.)}
\l{\cel{3}admisita por konkursar pri la Roma-Premio, dum ke li traktas}
\l{\cel{3}la temo propozita por la konkurso. - IV. Singla ek la chambreti}
\l{\cel{3}en qui la teatro-artisti chanjas sua pleo-vesti e fardizas su.}
\l{\cel{3}- V. Singla ek la mikra dividuri di spektaklo-chambrego, qua}
\l{\cel{3}kontenas tri, quar, kin o sis plasi. -VI. +Kluzajo (klozajo)}
\l{\cel{3}ube inter-asemblas grupo de framasoni, prezide da la veneracindo.}
\l{\cel{58}DEFR.}
\l{\sou{loko}. I. Parto determinita ek spaco. - II. Plaso, parto di spaco,}
\l{\cel{3}atribuita a kozo, a persono determinita. - III. Parto di la}
\l{\cel{3}texto di libro en qua esas texto determinita. - IV. Parto li-}
\l{\cel{3}mitizita en lando, en regiono. - DEFIRS.}
\l{}
\l{\sou{lokacar}\cel{2}(\sou{trans.)}\cel{1}Obtenar (de lua proprietanto) la yuro uzar, dum}
\l{\cel{3}periodo konvencionata, ulo po la pago di la preco konvencio-}
\l{\cel{3}nata. - F.}
\l{}
\l{\sou{lokativo.}\cel{1}(\sou{gram.)}\cel{1}Kazo gramatikala qua, en ula lingui, expresas}
\l{\cel{3}la loko, la destino. - DEFIS.}
\l{}
\l{loklo. Volvuro (en hari, rubando) forme di ringo. - DER.}
\l{}
\l{lokomobilo. Vapor-mashino transportebla. - DEFIRS.}
}
